\documentclass[11pt,a4paper]{article}

\usepackage[T1]{fontenc}
\usepackage[utf8]{inputenc}
\usepackage[french]{babel}
\usepackage{geometry}
\usepackage{enumitem}
\usepackage{hyperref}

\geometry{margin=2.3cm}
\setlist[itemize]{noitemsep, topsep=3pt}
\hypersetup{colorlinks=true, linkcolor=black, urlcolor=black}

\title{\vspace{-1.3cm}\textbf{Notes de pré-cadrage}\\[0.2cm]
\large FlowSense --- Challenge AIDGE 2026}
\author{EL BAROUD Elias \quad | \quad EL GAAMOUCH Aïssa \quad | \quad MAHI Rayan \quad | \quad MEBARKI Ilian}
\date{Septembre 2026}

\begin{document}
\maketitle
\thispagestyle{empty}

\section*{1. Problème à résoudre}

Dans des lieux isolés, notamment en montagne, il est utile de connaître l'occupation d'un espace sans installer une infrastructure lourde ni transmettre de vidéos vers le cloud. Le projet \textbf{FlowSense} vise à compter les personnes qui entrent et sortent d'une zone à travers une porte ou un couloir.

\section*{2. Objectif du projet}

Concevoir un capteur vidéo embarqué qui détecte les personnes et produit uniquement des données agrégées :

\begin{itemize}
    \item nombre d'entrées ;
    \item nombre de sorties ;
    \item estimation de l'occupation.
\end{itemize}

Le traitement doit être local, économe en énergie et compatible à terme avec le SoC \textit{ColibryNPU}. La cible du challenge est une consommation inférieure à \textbf{5 mW}.

\section*{3. Scénario de fonctionnement}

\begin{enumerate}
    \item Une caméra fixe observe une porte ou un couloir.
    \item Un modèle détecte les personnes dans chaque image.
    \item Un suivi associe les détections d'une même personne entre plusieurs images.
    \item La personne traverse une zone virtuelle définie à l'installation.
    \item Le système ajoute une entrée ou une sortie selon le sens de franchissement.
\end{enumerate}

La porte n'a pas besoin d'avoir de trait réel. Lors de l'installation, un utilisateur règle une fois la zone de passage dans l'image et indique quel côté correspond à l'intérieur. Cela permet d'utiliser le même système dans différents lieux, sans modifier le code.

\section*{4. Choix techniques provisoires}

\begin{itemize}
    \item \textbf{Prototype actuel :} YOLO nano pour détecter les personnes, puis ByteTrack pour leur attribuer un identifiant temporaire.
    \item \textbf{Calibrage :} deux limites virtuelles définissent une petite zone de porte ; un fichier JSON conserve ce réglage.
    \item \textbf{Comptage :} extérieur $\rightarrow$ intérieur = entrée ; intérieur $\rightarrow$ extérieur = sortie.
    \item \textbf{Version finale :} modèle beaucoup plus léger, entraîné avec des images variées, puis optimisé et quantifié en INT8 avec AIDGE avant déploiement sur ColibryNPU.
\end{itemize}

YOLO sert donc aujourd'hui à valider la chaîne \emph{détection $\rightarrow$ suivi $\rightarrow$ comptage}. Il ne constitue pas nécessairement le modèle final compatible avec la contrainte énergétique.

\section*{5. État actuel}

\begin{itemize}
    \item Première IA PyTorch \texttt{vide / personne} créée comme base d'apprentissage.
    \item Tests YOLO réalisés sur des images et des vidéos enregistrées par l'équipe.
    \item Premier comptage validé sur des vidéos de test : deux entrées séparées et une sortie individuelle.
    \item Outils en cours : calibrage de porte par clics et compteur vidéo configurable.
\end{itemize}

\section*{6. Prochaines étapes}

\begin{enumerate}
    \item Tester le compteur sur les scénarios \textit{porte vide}, \textit{une entrée}, \textit{une sortie}, \textit{deux personnes ensemble} et \textit{croisement}.
    \item Noter les erreurs : personne manquée, faux positif, identifiant perdu ou double comptage.
    \item Constituer un jeu de données varié : plusieurs portes, angles, éclairages et distances.
    \item Entraîner un détecteur plus léger et mesurer précision, rappel, latence et mémoire.
    \item Étudier la compatibilité AIDGE / ColibryNPU, la quantification INT8 et la consommation réelle.
\end{enumerate}

\section*{7. Questions à confirmer avec les organisateurs}

\begin{itemize}
    \item La limite de 5 mW inclut-elle la caméra, l'acquisition et le mode veille ?
    \item Comment la consommation énergétique sera-t-elle mesurée ?
    \item Quels opérateurs et quels formats de modèles sont pris en charge par ColibryNPU ?
    \item Quel scénario de comptage sera évalué : flux, distance, lumière, nombre de personnes ?
\end{itemize}

\vfill
\noindent\rule{\textwidth}{0.4pt}
\small Ces notes constituent un document de travail pour cadrer le projet avant la réception du matériel.

\end{document}
